% Generated draft: requires mathematical and source review.
\documentclass[11pt]{article}
\usepackage[margin=25mm]{geometry}
\usepackage{fontspec}
\setmainfont{Latin Modern Roman}
\usepackage{amsmath,amssymb,mathtools,unicode-math}
\setmathfont{Latin Modern Math}
\usepackage{enumitem,needspace,xurl,hyperref}
\setlength{\parindent}{0pt}
\setlength{\parskip}{5pt}
\setlength{\emergencystretch}{3em}
\newcommand{\N}{\mathbb{N}}
\newcommand{\Z}{\mathbb{Z}}
\newcommand{\Q}{\mathbb{Q}}
\newcommand{\R}{\mathbb{R}}
\newcommand{\C}{\mathbb{C}}
\newcommand{\F}{\mathbb{F}}
\newcommand{\A}{\mathbb{A}}
\newcommand{\PP}{\mathbb{P}}
\newcommand{\E}{\mathbb{E}}
\newcommand{\eps}{\varepsilon}
\newcommand{\dd}{\,\mathrm{d}}
\begin{document}
\section*{UnsolvedMath ID: 20001287}
\textbf{Source problem code:} catalogue entry
\subsection{Definitions and mathematical statement}
Let \(n\geq2\), let \(S^{n-1}\subseteq\R^n\) be the unit sphere, and let \(\sigma_{n-1}\) denote its surface measure. Retain the source's strict spherical dual
\[
 Q^\circ=\{x\in S^{n-1}:\langle x,y\rangle>0\ \text{for every }y\in Q\},\qquad
 P(Q)=\sigma_{n-1}(Q)\sigma_{n-1}(Q^\circ).
\]
A precise intended nondegenerate setting is \(Q=C\cap S^{n-1}\), where \(C\) is a full-dimensional pointed simplicial cone, generated by \(n\) linearly independent rays. Equivalently its bounding halfspaces have independent normals. The source says only ``determined by distinct \(n\) hyperplanes through the origin''; distinctness alone does not specify independence, chosen halfspaces or nondegeneracy.

Koenig's question, with the damaged inequality repaired according to the aid, is
\[
 P(Q)\leq P(\R_+^n\cap S^{n-1})
 =\frac{\sigma_{n-1}(S^{n-1})^2}{4^n}.
\]
The source states that the case \(n=3\) is true. No citation for that case is included.

For normalized notation, put \(\omega_d(C)=\Pr(g_d\in C)\), where \(g_d\) is a standard Gaussian in \(\R^d\), and let
\(C^*=\{x:\langle x,y\rangle\geq0\ \forall y\in C\}\).
For full-dimensional pointed polyhedral cones, strict and nonstrict dual boundaries have spherical measure zero, so
\[
 \frac{P(C\cap S^{n-1})}{\sigma_{n-1}(S^{n-1})^2}
 =\omega_n(C)\omega_n(C^*).
\]
Thus the conjectured normalized upper bound is \(4^{-n}\).

The supplied aid reports an exact structured family. Let \(W_\alpha\subseteq\R^2\) be a closed planar wedge of opening \(0<\alpha<\pi\), and put \(C_\alpha=W_\alpha\times\R_+^{n-2}\). For \(n=2\), the zero-dimensional factor has Gaussian angle one. It reports
\[
 \omega_n(C_\alpha)\omega_n(C_\alpha^*)
 =\frac{\alpha(\pi-\alpha)}{2^{2n-2}\pi^2}
 =4^{-n}-\frac{(\alpha-\pi/2)^2}{4^{n-1}\pi^2}.
\]
Hence the desired inequality holds throughout this family, with equality precisely at \(\alpha=\pi/2\), when the cone is an orthant up to orthogonal transformation. The deficit formula is exactly quadratic for this parameterized family; it is not a general stability theorem for all simplicial cones.

More generally the aid reports the orthogonal product rule. For cones \(C\subseteq\R^a\) and \(D\subseteq\R^b\),
\[
 (C\times D)^*=C^*\times D^*,\qquad
 \omega_{a+b}(C\times D)=\omega_a(C)\omega_b(D),
\]
and therefore their normalized volume products multiply. In particular, adjoining a \(b\)-dimensional positive orthant multiplies the product by \(4^{-b}\), lifting any verified lower-dimensional inequality to that orthant-product family. The aid reports this as partial progress, not a proof for arbitrary simplicial cones in all dimensions.
\subsection{Short English statement}
The question asks whether a spherical simplex and its strict dual have volume product at most that of an orthant. The aid reports an exact wedge-times-orthant formula, its equality case and deficit, plus a product rule; the general higher-dimensional question is left separate.
\subsection{Sources}
\begin{itemize}
\item Local registry: permanent UnsolvedMath ID \(20001287\), exact public code \texttt{catalogue entry}; \texttt{aim-convex-geometry-notes.json}, \texttt{notes[18]}.
\item AIM workshop ``Mahler's conjecture and duality in convex geometry'', attributed to H. Koenig: \url{https://aimath.org/WWN/mahlerduality/mahlerduality.pdf}. The source-item field and extraction print 3; the supplied verification identifies the intended entry as Problem 23. The source says the three-dimensional case is true. The literature field is empty and no authoritative separate citation for that assertion is supplied.
\item Research attempt 1: \texttt{partial\_result}, status at run \texttt{open}, rigor label \texttt{proved}, novelty confidence low, classification \texttt{PARTIAL-PROGRESS}. It reports the wedge--orthant formula and product-lifting principle as a candidate exact structured-family result. The current registry status is partially_solved. No separate publication bibliography is supplied.
\end{itemize}
\end{document}
